% GENERATED FILE. Edit canonical lessons, then run npm run generate:books.
% canonical-source-hash: fnv1a64:4e5ea041596aa0d3
% canonical-lessons: TE-C539-pallavi, TE-C539-caranam, TE-C539-padyam, TE-C539-geyam, TE-C539-janapadam

\chapter{Songs and Verse}
\label{ch:te-songs-and-verse}
\hlchaptermodality{539}

\begin{chapteropening}
\textbf{By the end of this chapter:} \emph{I can name the parts of a song and kinds of verse.}

Complete the last lesson of chapter 539: I can name the parts of a song and kinds of verse.
\end{chapteropening}

\section[\texorpdfstring{\te{పల్లవి} (pallavi)}{పల్లవి (pallavi)}]{\te{పల్లవి} (pallavi) — a refrain; a chorus}

\label{lesson:TE-C539-pallavi}



\begin{quote}
Before the new one: say the Telugu for a mridangam drum, then the Telugu for a reed pipe.
\end{quote}

\subsection*{You'll want to know: \te{పల్లవి}}
\textbf{\te{పల్లవి}} — \emph{pallavi} — \textquotedblleft{}a refrain; a chorus\textquotedblright{}.

\begin{quote}
\te{ఆ} \te{పాట} \te{పల్లవి} \te{అందరికీ} \te{గుర్తుంది}.
\end{quote}

\begin{quote}
\emph{ā pāṭa pallavi andarikī gurtundi.}
\end{quote}

\begin{quote}
Everyone remembers the refrain of that song.
\end{quote}

The opening line a song keeps returning to.

\begin{cousinweb}
\textbf{\te{పల్లవి}} (\emph{pallavi}) is usually traced to Sanskrit \emph{pallava}, \textquotedblleft{}a sprout, a twig, a blossom\textquotedblright{} (CDIAL 7969).
\end{cousinweb}

\begin{grammarlens}[title={What you've built}]
One more word for describing what you lived and felt.
\end{grammarlens}

\subsection*{Your turn}
\begin{itemize}
\raggedright
  \item \emph{Say it:} \emph{pallavi}
  \item \emph{Say it:} \emph{pallavi}, once more
  \item \emph{Say it:} \emph{ā pāṭa pallavi andarikī gurtundi.} — Everyone remembers the refrain of that song.
  \item \emph{Recall:} say the Telugu for a mridangam drum, then the Telugu for a reed pipe, then say \emph{pallavi} again
\end{itemize}

\begin{tcolorbox}[breakable,colback=teal!4,colframe=teal!35!black,title={Before you move on}]
What does \te{పల్లవి} mean? (\textquotedblleft{}A refrain; a chorus\textquotedblright{}.) Say it once more.
\end{tcolorbox}
\section[\texorpdfstring{\te{చరణం} (caraṇaṁ)}{చరణం (caraṇaṁ)}]{\te{చరణం} (caraṇaṁ) — a stanza; a verse of a song}

\label{lesson:TE-C539-caranam}



\begin{quote}
Before the new one: say the Telugu for a reed pipe, then the Telugu for a refrain.
\end{quote}

\subsection*{You'll want to know: \te{చరణం}}
\textbf{\te{చరణం}} — \emph{caraṇaṁ} — \textquotedblleft{}a stanza; a verse of a song\textquotedblright{}.

\begin{quote}
\te{రెండో} \te{చరణం} \te{నాకు} \te{గుర్తులేదు}.
\end{quote}

\begin{quote}
\emph{reṇḍō caraṇaṁ nāku gurtulēdu.}
\end{quote}

\begin{quote}
I do not remember the second verse.
\end{quote}

The verses that come between the refrains.

\begin{cousinweb}
\textbf{\te{చరణం}} (\emph{caraṇaṁ}) is borrowed from Sanskrit \emph{caraṇa}, \textquotedblleft{}a foot; a line of a stanza\textquotedblright{} (Monier-Williams).
\end{cousinweb}

\begin{grammarlens}[title={What you've built}]
One more word for describing what you lived and felt.
\end{grammarlens}

\subsection*{Your turn}
\begin{itemize}
\raggedright
  \item \emph{Say it:} \emph{caraṇaṁ}
  \item \emph{Say it:} \emph{caraṇaṁ}, once more
  \item \emph{Say it:} \emph{reṇḍō caraṇaṁ nāku gurtulēdu.} — I do not remember the second verse.
  \item \emph{Recall:} say the Telugu for a reed pipe, then the Telugu for a refrain, then say \emph{caraṇaṁ} again
\end{itemize}

\begin{tcolorbox}[breakable,colback=teal!4,colframe=teal!35!black,title={Before you move on}]
What does \te{చరణం} mean? (\textquotedblleft{}A stanza; a verse of a song\textquotedblright{}.) Say it once more.
\end{tcolorbox}
\section[\texorpdfstring{\te{పద్యం} (padyaṁ)}{పద్యం (padyaṁ)}]{\te{పద్యం} (padyaṁ) — a metrical verse; a poem}

\label{lesson:TE-C539-padyam}



\begin{quote}
Before the new one: say the Telugu for a refrain, then the Telugu for a stanza.
\end{quote}

\subsection*{You'll want to know: \te{పద్యం}}
\textbf{\te{పద్యం}} — \emph{padyaṁ} — \textquotedblleft{}a metrical verse; a poem\textquotedblright{}.

\begin{quote}
\te{తాతయ్య} \te{చాలా} \te{పద్యాలు} \te{నోటికి} \te{చెబుతారు}.
\end{quote}

\begin{quote}
\emph{tātayya cālā padyālu nōṭiki cebutāru.}
\end{quote}

\begin{quote}
Grandpa can recite many verses by heart.
\end{quote}

Plural: \te{పద్యాలు}, \emph{padyālu}.

Classical Telugu verse in a fixed metre; \te{పాట} is a song.

\begin{cousinweb}
\textbf{\te{పద్యం}} (\emph{padyaṁ}) is borrowed from Sanskrit \emph{padya}, \textquotedblleft{}relating to a foot; made of the parts of verses\textquotedblright{} (Monier-Williams).
\end{cousinweb}

\begin{grammarlens}[title={What you've built}]
One more word for describing what you lived and felt.
\end{grammarlens}

\subsection*{Your turn}
\begin{itemize}
\raggedright
  \item \emph{Say it:} \emph{padyaṁ}
  \item \emph{Say it:} \emph{padyaṁ}, once more
  \item \emph{Say it:} \emph{tātayya cālā padyālu nōṭiki cebutāru.} — Grandpa can recite many verses by heart.
  \item \emph{Recall:} say the Telugu for a refrain, then the Telugu for a stanza, then say \emph{padyaṁ} again
\end{itemize}

\begin{tcolorbox}[breakable,colback=teal!4,colframe=teal!35!black,title={Before you move on}]
What does \te{పద్యం} mean? (\textquotedblleft{}A metrical verse; a poem\textquotedblright{}.) Say it once more.
\end{tcolorbox}
\section[\texorpdfstring{\te{గేయం} (gēyaṁ)}{గేయం (gēyaṁ)}]{\te{గేయం} (gēyaṁ) — a lyric; a song text}

\label{lesson:TE-C539-geyam}



\begin{quote}
Before the new one: say the Telugu for a stanza, then the Telugu for a metrical verse.
\end{quote}

\subsection*{You'll want to know: \te{గేయం}}
\textbf{\te{గేయం}} — \emph{gēyaṁ} — \textquotedblleft{}a lyric; a song text\textquotedblright{}.

\begin{quote}
\te{ఆ} \te{కవి} \te{రాసిన} \te{గేయం} \te{బాగుంది}.
\end{quote}

\begin{quote}
\emph{ā kavi rāsina gēyaṁ bāgundi.}
\end{quote}

\begin{quote}
The lyric that poet wrote is good.
\end{quote}

Formal: verse written to be sung.

\begin{cousinweb}
\textbf{\te{గేయం}} (\emph{gēyaṁ}) is borrowed from Sanskrit \emph{gēya}, \textquotedblleft{}to be sung; singing, a song\textquotedblright{} (CDIAL 4249).
\end{cousinweb}

\begin{grammarlens}[title={What you've built}]
One more word for describing what you lived and felt.
\end{grammarlens}

\subsection*{Your turn}
\begin{itemize}
\raggedright
  \item \emph{Say it:} \emph{gēyaṁ}
  \item \emph{Say it:} \emph{gēyaṁ}, once more
  \item \emph{Say it:} \emph{ā kavi rāsina gēyaṁ bāgundi.} — The lyric that poet wrote is good.
  \item \emph{Recall:} say the Telugu for a stanza, then the Telugu for a metrical verse, then say \emph{gēyaṁ} again
\end{itemize}

\begin{tcolorbox}[breakable,colback=teal!4,colframe=teal!35!black,title={Before you move on}]
What does \te{గేయం} mean? (\textquotedblleft{}A lyric; a song text\textquotedblright{}.) Say it once more.
\end{tcolorbox}
\section[\texorpdfstring{\te{జానపదం} (jānapadaṁ)}{జానపదం (jānapadaṁ)}]{\te{జానపదం} (jānapadaṁ) — folk tradition; country life}

\label{lesson:TE-C539-janapadam}



\begin{quote}
Before the new one: say the Telugu for a metrical verse, then the Telugu for a lyric.
\end{quote}

\subsection*{You'll want to know: \te{జానపదం}}
\textbf{\te{జానపదం}} — \emph{jānapadaṁ} — \textquotedblleft{}folk tradition; country life\textquotedblright{}.

\begin{quote}
\te{జానపద} \te{పాటలు} \te{పల్లెల్లో} \te{పుట్టాయి}.
\end{quote}

\begin{quote}
\emph{jānapada pāṭalu pallellō puṭṭāyi.}
\end{quote}

\begin{quote}
Folk songs were born in the villages.
\end{quote}

Mostly as \te{జానపద}, folk: \te{జానపద} \te{కళలు}, folk arts.

\begin{cousinweb}
\textbf{\te{జానపదం}} (\emph{jānapadaṁ}) is borrowed from Sanskrit \emph{jānapada}, \textquotedblleft{}living in the country; a country dweller\textquotedblright{} (Monier-Williams).
\end{cousinweb}

\begin{grammarlens}[title={What you've built}]
One more word for describing what you lived and felt.
\end{grammarlens}

\subsection*{Your turn}
\begin{itemize}
\raggedright
  \item \emph{Say it:} \emph{jānapadaṁ}
  \item \emph{Say it:} \emph{jānapadaṁ}, once more
  \item \emph{Say it:} \emph{jānapada pāṭalu pallellō puṭṭāyi.} — Folk songs were born in the villages.
  \item \emph{Recall:} say the Telugu for a metrical verse, then the Telugu for a lyric, then say \emph{jānapadaṁ} again
\end{itemize}

\begin{tcolorbox}[breakable,colback=teal!4,colframe=teal!35!black,title={Before you move on}]
What does \te{జానపదం} mean? (\textquotedblleft{}Folk tradition; country life\textquotedblright{}.) Say it once more.
\end{tcolorbox}
